
    %%%%%%%%%%%%%%%%%%%%%%%%%%%%%%%%%%%%%%%%%%%%%%%%%%%%%%%%%%%%%%%%%%%%%%%%%%%%%%%%
    %%% Classe do documento
    \documentclass[12pt, addpoints]{exam}
    \UseRawInputEncoding

    %%%%%%%%%%%%%%%%%%%%%%%%%%%%%%%%%%%%%%%%%%%%%%%%%%%%%%%%%%%%%%%%%%%%%%%%%%%%%%%%
    % preambulo.tex
    %
    % Arquivo de configuração de packages para uso o LaTeX, conforme minhas
    % preferências e modelos pessoais.
    %
    % Para maiores informações, visite:
    %    https://github.com/abrantesasf/latex
    %
    % NÃO ALTERE SE NÃO SOUBER O QUE ESTÁ FAZENDO!
    %%%%%%%%%%%%%%%%%%%%%%%%%%%%%%%%%%%%%%%%%%%%%%%%%%%%%%%%%%%%%%%%%%%%%%%%%%%%%%%%


    %%%%%%%%%%%%%%%%%%%%%%%%%%%%%%%%%%%%%%%%%%%%%%%%%%%%%%%%%%%%%%%%%%%%%%%%%%%%%%%%
    %%% Estruturas de controle:
    \input{static/utils/controles}

    %%%%%%%%%%%%%%%%%%%%%%%%%%%%%%%%%%%%%%%%%%%%%%%%%%%%%%%%%%%%%%%%%%%%%%%%%%%%%%%%
    %%% Compilação condicional em PDF:
    \input{static/utils/pdf}

    %%%%%%%%%%%%%%%%%%%%%%%%%%%%%%%%%%%%%%%%%%%%%%%%%%%%%%%%%%%%%%%%%%%%%%%%%%%%%%%%
    %%% Configurações de layout da página:
    \input{static/utils/layout}

    %%%%%%%%%%%%%%%%%%%%%%%%%%%%%%%%%%%%%%%%%%%%%%%%%%%%%%%%%%%%%%%%%%%%%%%%%%%%%%%%
    %%% Configurações para autores:
    \input{static/utils/autores}

    %%%%%%%%%%%%%%%%%%%%%%%%%%%%%%%%%%%%%%%%%%%%%%%%%%%%%%%%%%%%%%%%%%%%%%%%%%%%%%%%
    %%% Cabeçalho e rodapé:
    %%% Atenção! É MUITO DIFÍCIL ajustar apropriadamente os running headers
    %%% e running footers da primeira vez. Você pode confiar no LaTeX e deixar que
    %%% ele faça o ajuste automaticamente ou pode quebrar a cabeça e
    %%% tentar melhorar algo que o LaTeX já faz muito bem, mesmo que não fique
    %%% exatamente do jeito que você quer. O que você prefere? Se quiser ajustar
    %%% manualmente, descomente a linha abaixo e faça as configurações no arquivo
    %%% "utils/headings.tex".
    %\input{static/utils/headings}

    %%%%%%%%%%%%%%%%%%%%%%%%%%%%%%%%%%%%%%%%%%%%%%%%%%%%%%%%%%%%%%%%%%%%%%%%%%%%%%%%
    %%% Configuração para as fontes do documento:
    %%% Se você quiser usar fontes próprias ou do sistema, precisará ajustar as
    %%% configurações no arquivo "utils/fontes.tex".
    %%% ATENÇÃO: por padrão eu utilizo XeLaTeX com fontes proprietárias projetadas
    %%% por Matthew Butterick, adquiridas em https://mbtype.com, que não podem ser
    %%% distribuídas devido à licença de utilização. Se você usar XeLaTeX, você
    %%% DEVERÁ OBRIGATORIAMENTE ajustar as fontes no arquivo "utils/fontes.tex".
    \input{static/utils/fontes}

    %%%%%%%%%%%%%%%%%%%%%%%%%%%%%%%%%%%%%%%%%%%%%%%%%%%%%%%%%%%%%%%%%%%%%%%%%%%%%%%%
    %%% Sumário:
    \input{static/utils/toc}

    %%%%%%%%%%%%%%%%%%%%%%%%%%%%%%%%%%%%%%%%%%%%%%%%%%%%%%%%%%%%%%%%%%%%%%%%%%%%%%%%
    %%% Matemática:
    \input{static/utils/matematica}

    %%%%%%%%%%%%%%%%%%%%%%%%%%%%%%%%%%%%%%%%%%%%%%%%%%%%%%%%%%%%%%%%%%%%%%%%%%%%%%%%
    %%% Notas de rodapé:
    \input{static/utils/footnotes}

    %%%%%%%%%%%%%%%%%%%%%%%%%%%%%%%%%%%%%%%%%%%%%%%%%%%%%%%%%%%%%%%%%%%%%%%%%%%%%%%%
    %%% Referências cruzadas, links, citações:
    \input{static/utils/crossrefs}

    %%%%%%%%%%%%%%%%%%%%%%%%%%%%%%%%%%%%%%%%%%%%%%%%%%%%%%%%%%%%%%%%%%%%%%%%%%%%%%%%
    %%% Configurações para os hiperlinks em PDF:
    \input{static/utils/pdfconfig}

    %%%%%%%%%%%%%%%%%%%%%%%%%%%%%%%%%%%%%%%%%%%%%%%%%%%%%%%%%%%%%%%%%%%%%%%%%%%%%%%%
    %%% Glossário e índice remissivo:
    \input{static/utils/glossindex}

    %%%%%%%%%%%%%%%%%%%%%%%%%%%%%%%%%%%%%%%%%%%%%%%%%%%%%%%%%%%%%%%%%%%%%%%%%%%%%%%%
    %%% Referências bibliográficas:
    \input{static/utils/bibliografia}

    %%%%%%%%%%%%%%%%%%%%%%%%%%%%%%%%%%%%%%%%%%%%%%%%%%%%%%%%%%%%%%%%%%%%%%%%%%%%%%%%
    %%% Cores:
    \input{static/utils/cores}

    %%%%%%%%%%%%%%%%%%%%%%%%%%%%%%%%%%%%%%%%%%%%%%%%%%%%%%%%%%%%%%%%%%%%%%%%%%%%%%%%
    %%% Computação:
    \input{static/utils/computacao}

    %%%%%%%%%%%%%%%%%%%%%%%%%%%%%%%%%%%%%%%%%%%%%%%%%%%%%%%%%%%%%%%%%%%%%%%%%%%%%%%%
    %%% Gráficos:
    \input{static/utils/graficos}

    %%%%%%%%%%%%%%%%%%%%%%%%%%%%%%%%%%%%%%%%%%%%%%%%%%%%%%%%%%%%%%%%%%%%%%%%%%%%%%%%
    %%% Ícones:
    \input{static/utils/icones}

    %%%%%%%%%%%%%%%%%%%%%%%%%%%%%%%%%%%%%%%%%%%%%%%%%%%%%%%%%%%%%%%%%%%%%%%%%%%%%%%%
    %%% Blocos:
    \input{static/utils/blocos}

    %%%%%%%%%%%%%%%%%%%%%%%%%%%%%%%%%%%%%%%%%%%%%%%%%%%%%%%%%%%%%%%%%%%%%%%%%%%%%%%%
    %%% Boxes coloridos:
    \input{static/utils/colorbox}

    %%%%%%%%%%%%%%%%%%%%%%%%%%%%%%%%%%%%%%%%%%%%%%%%%%%%%%%%%%%%%%%%%%%%%%%%%%%%%%%%
    %%% Tabelas:
    \input{static/utils/tabelas}

    %%%%%%%%%%%%%%%%%%%%%%%%%%%%%%%%%%%%%%%%%%%%%%%%%%%%%%%%%%%%%%%%%%%%%%%%%%%%%%%%
    %%% Ambientes de listas:
    \input{static/utils/listas}

    %%%%%%%%%%%%%%%%%%%%%%%%%%%%%%%%%%%%%%%%%%%%%%%%%%%%%%%%%%%%%%%%%%%%%%%%%%%%%%%%
    %%% Ambientes floats e similares:
    \input{static/utils/floats}

    %%%%%%%%%%%%%%%%%%%%%%%%%%%%%%%%%%%%%%%%%%%%%%%%%%%%%%%%%%%%%%%%%%%%%%%%%%%%%%%%
    %%% Ajustes para a classe EXAM:
    \input{static/utils/exam}

    %%%%%%%%%%%%%%%%%%%%%%%%%%%%%%%%%%%%%%%%%%%%%%%%%%%%%%%%%%%%%%%%%%%%%%%%%%%%%%%%
    %%% Ajustes para textos:
    \input{static/utils/texto}

    %%%%%%%%%%%%%%%%%%%%%%%%%%%%%%%%%%%%%%%%%%%%%%%%%%%%%%%%%%%%%%%%%%%%%%%%%%%%%%%%
    %%% Ambientes, aliases, comandos e customizações em geral para serem
    %%% utilizadas nos documentos. Defina aqui qualquer customização específica
    %%% para seus documentos!
    \input{static/utils/customizacoes}

    %%%%%%%%%%%%%%%%%%%%%%%%%%%%%%%%%%%%%%%%%%%%%%%%%%%%%%%%%%%%%%%%%%%%%%%%%%%%%%%%
    %%% Ajuste do layout, espaçamento de linhas e etc.:
    \geometry{a4paper, portrait, left=2.0cm, right=2.0cm, top=2.0cm, bottom=2.0cm}
    \singlespacing

    %%%%%%%%%%%%%%%%%%%%%%%%%%%%%%%%%%%%%%%%%%%%%%%%%%%%%%%%%%%%%%%%%%%%%%%%%%%%%%%%
    %%% Configurações de cabeçalho e rodapé (não use fancyhdr pois ocorre conflito):
    \pagestyle{headandfoot}
    \runningheadrule
    \firstpageheader{}{}{}
    \runningheader{Sem disciplina}{}{Avaliação}
    \firstpagefooter{}{}{Página \thepage\ de \numpages}
    \runningfooter{}{}{Página \thepage\ de \numpages}
    \runningfootrule

    %%%%%%%%%%%%%%%%%%%%%%%%%%%%%%%%%%%%%%%%%%%%%%%%%%%%%%%%%%%%%%%%%%%%%%%%%%%%%%%%
    %%% Configurações para as propriedades do PDF:
    \hypersetup{
    hidelinks,           % Comente para web, descomente para impressão
    colorlinks=true,      % True para web, False para impressão
    pdftitle=BassulMolestado: Avaliação,
    pdfauthor=Matheus Endlich Silveira,
    pdfsubject=Sem disciplina,
    pdfkeywords=['Sem tag'],
    pdfinfo={
        CreationDate={}, % Ex.: D:AAAAMMDDHH24MISS
        ModDate={}       % Ex.: D:AAAAMMDDHH24MISS
    }
    }
    \input{static/utils/pdfconfig}

    %%%%%%%%%%%%%%%%%%%%%%%%%%%%%%%%%%%%%%%%%%%%%%%%%%%%%%%%%%%%%%%%%%%%%%%%%%%%%%%%
    %%% Começa o documento
    \begin{document}

    %%%%%%%%%%%%%%%%%%%%%%%%%%%%%%%%%%%%%%%%%%%%%%%%%%%%%%%%%%%%%%%%%%%%%%%%%%%%%%%%
    %%% Folha de instruções padronizadas da UVV para a prova
    \pdfbookmark[1]{Instruções}{instrucoes}
    \begin{figure}[H]
        \begin{center}
            \includegraphics[scale=0.3]{{static/imgs/logo-uvv.png}}
        \end{center}	
    \end{figure}

    \vspace{-1cm}

    \begin{table}[H]
        \centering
        \begin{tabular}{|llll|}
            \hline
            \multicolumn{2}{|l|}{\textbf{Disciplina:} Sem disciplina} & 
            \multicolumn{1}{l|}{\multirow{3}{*}{\textbf{\makecell{Nota:\\ \,\\ \,}}}} & 
            \multirow{3}{*}{\textbf{\makecell{Coordenador:\\ \,\\ \,}}} \\ 
            \cline{1-2}
            \multicolumn{2}{|l|}{\textbf{Professor:} Matheus Endlich Silveira \hspace{4.72cm} } & 
            \multicolumn{1}{l|}{} & \\ 
            \cline{1-2}
            \multicolumn{2}{|l|}{\textbf{Aluno:}} & 
            \multicolumn{1}{l|}{} & \\ 
            \hline
            \multicolumn{1}{|l|}{\textbf{Turma:} \hspace{6.3cm} } & 
            \multicolumn{1}{l|}{\textbf{Semestre:}} & 
            \multicolumn{2}{l|}{\textbf{Valor:} 10 pontos} \\ 
            \hline
            \multicolumn{1}{|l|}{\textbf{Data:}} & 
            \multicolumn{3}{l|}{\textbf{Avaliação:}} \\ 
            \hline
        \end{tabular}
    \end{table}

    \begin{center}
        \fbox{\setlength\fboxsep{0.3cm} \fbox{\parbox{15.4cm}{
            \begin{center}
                \textbf{INSTRUÇÕES PARA A REALIZAÇÃO DESTA AVALIAÇÃO:}
            \end{center}
            \begin{itemize}[leftmargin=*]
                \item Esta é a primeira avaliação da disciplina Sem disciplina, 
                    e engloba o conteúdo referente Sem tag.
                \item Esta avaliação contém 50 questões objetivas, \textbf{todas obrigatórias}.
                \item \textbf{Leia atentamente cada questão!} As questões objetivas terão uma,
                    e apenas uma única, resposta a ser marcada.
                \item \textbf{Desligue o celular} antes de começar. A avaliação é
                    \textbf{individual} e \textbf{sem consulta}.
                \item As questões podem ser respondidas, na prova, com lápis ou caneta. Ao final
                    da prova \textbf{{VOCÊ DEVERÁ PREENCHER O GABARITO COM CANETA
                    DE COR PRETA}} \textbf{\underline{OU AZUL}} para que seja feita a correção
                    automatizada através da leitura do padrão de respostas marcado no
                    gabarito.
                \item \textbf{CUIDADO ao preencher o gabarito} pois eles são \textbf{INDIVIDUAIS
                    e NOMEADOS} (cada estudante tem um gabarito próprio específico, com um
                    código de identificação único). \textbf{Questões rasuradas são anuladas}
                    pelo software de leitura do gabarito.
                \item Ao final da prova, \textbf{DEVOLVA A PROVA E O GABARITO} para o professor.
                \item Siga todas as normas de \textbf{Integridade Acadêmica} da disciplina.
                    Alunos que forem flagrados com qualquer espécie de ``cola'' ou trocando
                    informações com outros alunos terão suas avaliações recolhidas, as notas
                    zeradas, e a situação será encaminhada para a coordenação para a aplicação
                    das penalidades previstas pela universidade.
                \item Com 90 minutos de avaliação você terá tempo de até 2,min por questão,
                    em média.
                \item Boa avaliação!
            \end{itemize}
            \vspace{2cm}
        }}}
    \end{center}
    \newpage

    %%%%%%%%%%%%%%%%%%%%%%%%%%%%%%%%%%%%%%%%%%%%%%%%%%%%%%%%%%%%%%%%%%%%%%%%%%%%%%%%
    %%% Inicia o bloco de questões:
    \begin{questions}

    \newpage
    \question

Dentre as definições a seguir, ligadas ao conceito de visões do modelo relacional, qual delas é \textbf{INCORRETA}?


\begin{choices}
\choice Uma visão é útil por representar uma percepção particular do banco de dados, compartilhado por muitos aplicativos.
\choice O gerenciamento de visões envolve a conversão da consulta do usuário sobre as visões para a consulta sobre as relações base.
\choice Programas aplicativos do banco de dados podem ser executados sobre visões de relações da base de dados.
\choice Uma visão relacional é uma relação virtual que nunca é materializada.
\choice Uma visão relacional é uma relação virtual, derivada de relações base a partir da especificação de operações da álgebra relacional.
\end{choices}

\question

Uma definição lógica, abstrata e auto-contida dos objetos que modelam a

estrutura de armazenamento de dados, dos operadores que modelam o comportamento

e dos demais conceitos que modelam a integridade, e que, juntos, constituem a

máquina abstrata com a qual os usuários interagem é chamada de:


\begin{choices}
\choice Banco de dados
\choice Nenhuma das respostas acima
\choice Modelo de dados
\choice Modelo entidade-relacionamento
\choice Modelo relacional
\end{choices}

\question

A arquitetura ilustrada na figura abaixo é um exemplo típico de qual sistema?

(CU = unidade de controle; IS = \ingles{stream} de instruções; PU = unidade de

processamento; DS = \ingles{stream} de dados; LM = memória local)

\begin{figure}[H]

\begin{center}

\includegraphics[width=0.5\linewidth]{static/imgs/defaul.png}

\end{center}

\end{figure}

\vspace{-1cm}
\begin{choices}
\choice Multiprocessadores simétricos
\choice Acesso não uniforme à memória
\choice Processadores multicore
\choice Clusters
\choice Processadores seqüenciais
\end{choices}

\question

A média aritmética de uma lista de 50 números é 50. Se dois desses números, 51 e 97, forem suprimidos dessa lista a média dos restantes será


\begin{choices}
\choice 47
\choice 49
\choice 50
\choice 51
\choice 40
\end{choices}

\question

\textit{Starvation} ocorre quando:
\begin{choices}
\choice Pelo menos um processo é continuamente postergado e não executa.
\choice Dois ou mais processos são forçados a acessar dados críticos alternando estritamente entre eles.
\choice Pelo menos um evento espera por um evento que não vai ocorrer.
\choice O processo tenta mas não consegue acessar uma variável compartilhada.
\choice A prioridade de um processo é ajustada de acordo com o tempo total de execução do mesmo.
\end{choices}

\question

Seu chefe que fazer uma campanha de venda para produtos que custam entre 20 e 30

reais, mas está preocupado porque ouviu de um dos vendedores que diversos

produtos nessa faixa de preço estão com o estoque zerado. Para verificar a

situação seu chefe pediu para que você prepare um relatório que mostre os

produtos que custam entre 20 e 30 reais e que estão com o estoque zerado. Você

preparou o seguinte relatório para seu chefe:

\begin{tcolorbox}

 \begin{verbatim}1  SELECT l.nome AS loja

2       , p.nome AS produto

3       , p.preco_unitario AS preco

4       , e.quantidade AS qtd

5  FROM produtos p

6  INNER JOIN estoques e ON (e.produto_id = p.produto_id)

7  INNER JOIN lojas l ON (e.loja_id = l.loja_id)

8  WHERE (p.preco_unitario >= 20 OR p.preco_unitario <= 30)

9    AND e.quantidade = 0

10 ORDER BY loja, produto;

 \end{verbatim}

\end{tcolorbox}

O seu relatório:
\begin{choices}
\choice Está errado pois a tabela que deveria estar na cláusula FROM é a tabela ``lojas', não a tabela ``produtos'.
\choice Nenhuma das alternativas anteriores.
\choice Está errado pois há um erro de sintaxe nas linhas 6 e 7.
\choice Está errado pois vai trazer também produtos fora da faixa de preço que o seu chefe quer.
\choice Está correto e trará exatamente, os produtos com o estoque zerado e na faixa de preço requisitada pelo seu chefe. Além disso ordena o resultado pelo nome da loja e pelo nome do produto.
\end{choices}

\question

O conjunto básico de atividades e a ordem em que são realizadas no processo de construção de um software definem o que é habitualmente denominado de ciclo de vida do software. O ciclo de vida tradicional (também denominado waterfall) ainda é hoje em dia um dos mais difundidos e tem por característica principal: 
\begin{choices}
\choice o uso de formalização rigorosa em todas as etapas de desenvolvimento; 
\choice a abordagem sistemática para realização das atividades do desenvolvimento de software de modo que elas seguem um fluxo sequencial; 
\choice a codificação de uma versão executável do sistema desde as fases iniciais do desenvolvimento, de modo que o sistema final é incrementalmente construído, daí a alusão à ideia de "cascata" (waterfall); 
\choice a priorização da análise dos riscos do desenvolvimento
\choice a avaliação constante dos resultados intermediários feita pelo cliente.
\end{choices}

\question

O processo de visualização de objetos 3D envolve uma série de passos desde a representação vetorial de um objeto até a exibição da imagem correspondente na tela do computador (pipeline 3D). Selecione a alternativa abaixo que reflete a ordem correta em que esses passos devem ocorrer. 
\begin{choices}
\choice Transformação de câmera, recorte 3D, projeção, mapeamento para coordenadas de tela, rasterização. 
\choice Transformação de câmera, mapeamento para coordenadas de tela, recorte 3D, rasterização, projeção. 
\choice Recorte 3D, transformação de câmera, rasterização, projeção, mapeamento para coordenadas de tela. 
\choice Projeção, transformação de câmera, recorte 3D, mapeamento para coordenadas de tela, rasterização. 
\choice Nenhuma das respostas acima está correta.
\end{choices}

\question

``Um modelo de dados é uma definição lógica, abstrata e auto-contida

dos \circled{1}, \circled{2} e \circled{3} que, juntos, constituem a \circled{4}

com a qual os usuários interagem. Os \circled{1} nos permitem modelar

a \circled{5}; os \circled{2} nos permitem modelar

seu \circled{6}. Os \circled{3} nos permitem modelar sua \circled{7}.'



As palavras/termos que completam a frase acima, são, conforme os números:


\begin{choices}
\choice 1 = operadores; 2 = objetos; 3 = demais conceitos; 4 = máquina abstrata; 5 = estrutura de dados; 6 = comportamento; 7 = integridade.
\choice 1 = demais conceitos; 2 = objetos; 3 = operadores; 4 = estrutura de dados; 5 = máquina abstrata; 6 = integridade; 7 = comportamento.
\choice 1 = objetos; 2 = operadores; 3 = máquina abstrata; 4 = demais conceitos; 5 = estrutura de dados; 6 = integridade; 7 = comportamento.
\choice 1 = objetos; 2 = operadores; 3 = demais conceitos; 4 = máquina abstrata; 5 = estrutura de dados; 6 = comportamento; 7 = integridade.
\choice 1 = objetos; 2 = máquina abstrata; 3 = demais conceitos; 4 = operadores; 5 = estrutura de dados; 6 = comportamento; 7 = integridade.
\end{choices}

\question

O menor número possível de arestas em um grafo conexo com \( n \) vértices é:
\begin{choices}
\choice \( n - 1 \)
\choice \( n^2 \)
\choice \( n/2 \)
\choice \( n \)
\choice 1
\end{choices}

\question

Suponha que \( T \) seja uma árvore AVL inicialmente vazia, e considere a inserção dos elementos \(10, 20, 30, 5, 15, 2\) em \( T \), nesta ordem. Qual das sequências abaixo corresponde a um percurso de \( T \) em pré-ordem:
\begin{choices}
\choice 30, 20, 15, 10, 5, 2
\choice 15, 10, 5, 2, 20, 30
\choice 2, 5, 10, 15, 20, 30
\choice 20, 10, 5, 2, 15, 30
\choice 10, 5, 2, 20, 15, 30
\end{choices}

\question

Pode-se afirmar que o gráfico da função \(y = 2 + \frac{1}{x - 1}\) é o gráfico da função \(y = \frac{1}{x}\):
\begin{choices}
\choice nenhuma das anteriores.
\choice transladado uma unidade para a direita e duas unidades para cima
\choice transladado uma unidade para a direita e duas unidades para baixo
\choice transladado uma unidade para a esquerda e duas unidades para baixo
\choice transladado uma unidade para a esquerda e duas unidades para cima
\end{choices}

\question

Considere a seguinte tabela para uma base de dados relacional:\\

Empregado (CodEmp, NomeEmp, CodDepto)\\

Considere que esta tabela tem um índice na forma de uma árvore B sobre as colunas (CodEmp, CodDepto), nesta ordem.

Quanto a este índice, considere as seguintes afirmativas:

\begin{enumerate}

    \item Este índice pode ser usado pelo SGBD relacional para acelerar uma consulta na qual são fornecidos os valores de CodEmp e CodDepto.

    \item Este índice pode ser usado pelo SGBD relacional para acelerar uma consulta na qual é fornecido um valor de CodEmp.

    \item Este índice não é adequado para ser usado pelo SGBD relacional para acelerar uma consulta na qual é fornecido um valor de CodDepto.

    \item O algoritmo que faz inserções e remoções de entradas do índice tem por objetivo garantir que o índice fique organizado de tal forma que o acesso a cada nodo da árvore implique em número de acessos semelhantes.

    \item O índice por árvore-B não é adequado para tabelas que sofrem grande número de inclusões e exclusões, pois exige reorganizações frequentes.

\end{enumerate}





Quanto a estas afirmativas pode se dizer que:
\begin{choices}
\choice Apenas as afirmativas 1), 2) e 5) estão corretas
\choice Apenas as afirmativas 1), 2), 3) e 4) estão corretas
\choice Apenas as afirmativas 1), 2) e 4) estão corretas
\choice Todas afirmativas estão corretas
\choice Nenhuma das afirmativas está correta
\end{choices}

\question

São vantagens da utilização de threads no espaço do usuário, exceto:
\begin{choices}
\choice O escalonamento pode ser específico para a aplicação.
\choice Maior portabilidade da aplicação.
\choice A criação e o gerenciamento das threads são mais eficientes.
\choice O sistema operacional escalona a thread.
\choice Nenhuma modificação é necessária no kernel.
\end{choices}

\question

Ao executar o seguinte comando SQL, nenhum dado foi retornado, ou seja, não

ocorreu nenhum erro de sintaxe mas nenhum dado foi encontrado:

\begin{tcolorbox}

 \begin{verbatim}1   SELECT c.nome AS cliente

 2   FROM clientes c

 3   LEFT JOIN pedidos p ON (p.cliente_id = c.cliente_id)

 4   WHERE p.cliente_id IS NULL

 5   ORDER BY cliente;

 \end{verbatim}

\end{tcolorbox}

Por que esse comando SQL não retornou nenhum dado?


\begin{choices}
\choice Há um erro lógico na linha 4.
\choice A junção que deveria ter sido realizada, na linha 3, era uma RIGHT JOIN ao invés de uma LEFT JOIN.
\choice Há um erro semântico na linha 3
\choice Todos os clientes cadastrados fizeram, pelo menos, um pedido de compra.
\choice A junção que deveria ter sido realizada, na linha 3, era uma FULL JOIN ao invés de uma RIGHT JOIN.
\end{choices}

\question

Assinale a proposição verdadeira


\begin{choices}
\choice Se \( x \) é um número real tal que \( x^2 \leq 4 \) então \( x \leq 2 \) e \( x \leq -2 \)
\choice Se \( x < -4 \) ou \( x > 1 \) então \( \frac{2x + 3}{x - 1} > 1 \)
\choice nenhuma das anteriores
\choice Se \( x + y \) é um número racional então \( x \) e \( y \) são números racionais
\choice Se \( x \) e \( y \) são números reais tais que \( x < y \) então \( x^2 < y^2 \)
\end{choices}

\question

Qual o significado de coerência de memórias \textit{cache} em sistemas multiprocessados?


\begin{choices}
\choice Caches em processadores diferentes sempre contêm o mesmo dado válido para a mesma linha de cache.
\choice Caches em processadores diferentes nunca interagem entre si.
\choice Caches em processadores diferentes sempre lêem os mesmos dados ao mesmo tempo.
\choice Caches em processadores diferentes podem possuir dados diferentes associados à mesma linha de cache.
\choice Caches em processadores diferentes nunca compartilham a mesma linha de cache.
\end{choices}

\question

Considere o circuito abaixo, implementado com duas portas \textbf{NAND}.



\begin{figure}[H]

    \centering

    \includegraphics[width=0.5\linewidth]{static/imgs/defaul.png}

    \caption{Questão 23}

    \label{fig:enter-label}

\end{figure}



Qual das seguintes portas equivale a este circuito?
\begin{choices}
\choice \textbf{AND}
\choice \textbf{XOR}
\choice \textbf{OR}
\choice \textbf{NOT}
\choice \textbf{NOR}
\end{choices}

\question

Qual das afirmações a seguir, relativas à análise sintática, está INCORRETA?


\begin{choices}
\choice Uma das diferenças entre os diversos algoritmos de análise redutiva é a forma de identificar o \textit{handle} na pilha.
\choice Analisadores sintáticos descendentes recursivos são mais simples de implementar do que analisadores sintáticos redutivos.
\choice As gramáticas LL podem descrever mais linguagens do que as gramáticas LR.
\choice Algoritmos descendentes recursivos podem ser utilizados para algumas gramáticas ambíguas.
\choice Algoritmos de análise redutiva podem ser utilizados mesmo para gramáticas ambíguas.
\end{choices}

\question

Considere o seguinte trecho de programa:\\

1. \( i := 1; \)\\

2. while \( i \leq n \) do\\

   begin\\

3. \( sum := sum + a[i]; \)\\

4. \( i := i + 1; \)\\

   end;\\



Considere que:\\

- \( I \) representa a inicialização da variável \( i := 1 \) na linha 1;\\

- \( T \) representa o teste da linha 2;\\

- \( A \) representa os comandos da linha 3;\\

- \( P \) representa o incremento na linha 4.\\



Qual é a expressão regular que representa todas as sequências de passos possíveis de serem executados por este trecho de programa?
\begin{choices}
\choice \( IT^+A^*P^* \)
\choice \( I(TAP)^* \)
\choice \( I(TAP)^+ \)
\choice \( I(T(AP)^*)^+ \)
\choice \( I(T(AP)^*)^* \)
\end{choices}

\question

Qual é a função do circuito da figura abaixo?

\begin{figure}

    \centering

    \includegraphics[width=0.5\linewidth]{static/imgs/defaul.png}

    \label{fig:enter-label}

\end{figure}
\begin{choices}
\choice multiplexador
\choice somador
\choice subtrator
\choice deslocador
\end{choices}

\question

O determinante da matriz dada abaixo é

$\begin{pmatrix}

2 & 7 & 9 & -1 & 1 \\

2 & 8 & 3 & 1 & 0 \\

-1 & 0 & 4 & 3 & 0 \\

2 & 0 & -1 & 0 & 0 \\

3 & 0 & 0 & 0 & 0 \\

\end{pmatrix}$
\begin{choices}
\choice -96
\choice 96
\choice 86
\choice 46
\choice -86
\end{choices}

\question

Histograma de uma imagem com \( K \) tons de cinza é:
\begin{choices}
\choice Contagem do número de tons de cinza que ocorreram na imagem.
\choice Contagem do número de objetos encontrados na imagem.
\choice Nenhuma alternativa acima.
\choice Contagem dos pixels da imagem.
\choice Contagem do número de vezes que cada um dos \( K \) tons de cinza ocorreu na imagem.
\end{choices}

\question

Se \( \lim_{n \to \infty} \left( \frac{1}{n^2} + \frac{2}{n^2} + \cdots + \frac{n-1}{n^2} \right) = L \) então


\begin{choices}
\choice \( L = 0 \)
\choice \( L = 1/2 \)
\choice \( L = \infty \)
\choice \( L = 1 \)
\choice \( L = 2 \)
\end{choices}

\question

Considere as seguintes tabelas em uma base de dados relacional:\\

\textbf{Departamento (CodDepto, NomeDepto)}\\

\textbf{Empregado (CodEmp, NomeEmp, CodDepto)}\\

Deseja-se obter uma tabela na qual cada linha é a concatenação de uma linha da tabela Departamento com uma linha da tabela de Empregado. Caso um departamento não possua empregados, sua linha no resultado deve conter vazio (NULL) nos campos referentes ao empregado. A operação de álgebra relacional que deve ser aplicada para combinar estas duas tabelas é:
\begin{choices}
\choice Junção externa
\choice Junção interna
\choice Projeção
\choice União
\choice Divisão
\end{choices}

\question

São linguagens clássicas que representam aplicações científicas, aplicações 

comerciais, aplicações de inteligência artificial, e aplicações gerais,

respectivamente:
\begin{choices}
\choice COBOL, Java, LISP, Fortran
\choice R, COBOL, C, Python
\choice Fortran, COBOL, LISP, C
\choice MATLAB, COBOL, SML, LISP
\choice Java, MATLAB, Scheme, JavaScript
\end{choices}

\question

Dada a seguinte fórmula (lógica de primeira ordem):\\

$\forall x \ \exists y \ | \ ama(x,y)$ \\

qual das seguintes sentenças em linguagem natural ela representa, considerando que \( ama(x,y) \) representa que \( x \) ama \( y \)?


\begin{choices}
\choice Há alguém que todos amam.
\choice Nenhuma das anteriores.
\choice Ninguém ama a todos.
\choice Todos amam alguém.
\choice Alguém ama a todos.
\end{choices}

\question

Considere as seguintes tabelas em uma base de dados relacional:\\

Departamento (\underline{CodDepto}, NomeDepto)\\

Empregado (\underline{CodEmp}, NomeEmp, CodDepto, Salario)\\

Considere a seguinte consulta SQL:

\begin{verbatim}

SELECT d.CodDepto, d.NomeDepto, SUM(e.Salario)

FROM Departamento d

JOIN Empregado e ON d.CodDepto = e.CodDepto

GROUP BY d.CodDepto, d.NomeDepto

HAVING COUNT(e.CodEmp) > 2 AND AVG(e.Salario) > 40;

\end{verbatim}



Qual o resultado da consulta?
\begin{choices}
\choice Para cada departamento que tem mais que dois empregados e cuja média salarial é maior que 40, obter um grupo de linhas que contém, para cada empregado do departamento, o código de seu departamento, seguido do nome de seu departamento, seguido da soma dos salários dos empregados do departamento.
\choice A consulta não retorna nada pois está incorreta.
\choice Para cada departamento que tem mais que dois empregados e cuja média salarial é maior que 40, obter o código de departamento, seguido do nome do departamento, seguido da soma dos salários dos empregados do departamento.
\choice Para cada empregado que tem mais que dois departamentos, ambos com média salarial maior que 40, obter o código de departamento, seguido do nome do departamento, seguido da soma dos salários dos empregados do departamento.
\choice Para cada departamento que tem mais que dois empregados e cuja média salarial, considerando todos empregados do departamento, exceto os dois primeiros, é maior que 40, obter o código de departamento, seguido do nome do departamento, seguido da soma dos salários dos empregados do departamento.
\end{choices}

\question

A velocidade de um ponto em movimento é dada pela equação\\

$v(t) = t e^{-0.01t} \, \text{m/s}$\\

O espaço percorrido desde o instante que o ponto começou a se mover até a sua parada total é


\begin{choices}
\choice \( 10^2 \, m \)
\choice \( 10^3 e^{-0.01} \, m \)
\choice \( 10^4 \, m \)
\choice \( 10^2 e^{-1} \, m \)
\choice \( (e^{-100} - 1) m \)
\end{choices}

\question

Quando Edgar Codd cricou o modelo relacional de dados,

também criou algumas regras para que a estrutura das tabelas fosse adequada às

operações de inserção, atualização, remoção e busca de dados. Ao avaliar um

banco de dados que armazena dados dos eleitores, dos partidos e dos candidatos

de uma eleição, você encontrou a situação ilustrada na tabela abaixo:

\begin{figure}[H]

  \centering

  \includegraphics[width=0.5\linewidth]{static/imgs/defaul.png}

\end{figure}

Em relação às boas normas de projeto do modelo relacional, assinale a

alternativa que corretamente identifica um problema no projeto desta tabela:
\begin{choices}
\choice Todas as afirmativas anteriores.
\choice Apesar do telefone ser um atributo multivalorado, do jeito que a tabela está criada só é possível armazenar dois números de telefones, criando uma restrição ao armazenamento de telefones de contado de eleitores que tenham mais de dois telefones.
\choice A tabela está misturando dados de três entidades diferentes, eleitores, partidos e candidatos (o correto seria que cada tabela armazenasse dados de um único assunto ou entidade).
\choice A tabela contém um atributo multivalorado (o correto seria que cada tabela não tivesse atributos multivalorados).
\choice A chave primária para a tabela não está definida e, somente olhando as colunas, não é possível dizer se é a coluna 'numero' ou a coluna 'titulo'.
\end{choices}

\question

Os protocolos de transporte atribuem a cada serviço um identificador único, o qual é empregado para encaminhar uma requisição de um aplicativo cliente ao processo servidor correto. Nos protocolos de transporte TCP e UDP, como esse identificador se denomina?
\begin{choices}
\choice Protocolo de aplicação
\choice Endereço IP
\choice Porta
\choice Identificador do processo (PID)
\choice Conexão
\end{choices}

\question

Em um banco de dados relacional, o que é uma chave primária?


\begin{choices}
\choice É qualquer coluna ou conjunto de colunas que podem servir como índice
\choice Nenhuma das respostas acima
\choice É denominada por FK
\choice É um atributo (ou um conjunto de atributos) que permite identificar única e exclusivamente cada registro da tabela
\choice É uma validação que evita o cadastro de dados errados
\end{choices}

\question

Quais das seguintes afirmações são verdadeiras? As Métricas de software servem para:



\begin{itemize}

    \item[I)] indicar a qualidade do produto e avaliar a produtividade.

    \item[II)] auxiliar na melhoria do processo.

    \item[III)] formar uma base para as estimativas e justificar a aquisição de ferramentas.

    \item[IV)] determinar se a utilização de um método traz benefícios ou não.

\end{itemize}
\begin{choices}
\choice Apenas as alternativas I, II e IV.
\choice Todas as alternativas.
\choice Apenas as alternativas I, IV.
\choice Nenhuma delas.
\choice Apenas as alternativas II e III.
\end{choices}

\question

Sobre a UML, quais das seguintes afirmações são verdadeiras?



\begin{itemize}

    \item[I)] A UML é o método de desenvolvimento de software mais utilizado na atualidade.

    \item[II)] A UML é uma evolução das linguagens para especificação dos conceitos dos métodos de Booch, OMT e OOSE e também de outros métodos de especificação de requisitos de software orientados a objetos ou não.

    \item[III)] A UML é composta dos seguintes diagramas: Diagrama de Caso de Uso, Diagrama de Classes, Diagrama de Colaboração, Diagrama de Estados, entre outros.

    \item[IV)] Em UML pode-se representar tão somente relacionamentos de Agregação, Associação e Composição.

\end{itemize}
\begin{choices}
\choice Apenas as alternativas II e III.
\choice Apenas as alternativas I, II e III.
\choice Apenas as alternativas III e IV.
\choice Nenhuma delas.
\choice Todas as alternativas.
\end{choices}

\question

\textbf Supondo a Relação \texttt{PROJ (PNO, Nome, Orçam)}, com chave primária \texttt{PNO} e a Relação \texttt{DSG (ENO, PNO, Dur, Resp)}, com chave primária \{\texttt{ENO, PNO}\} e chave estrangeira \texttt{PNO} em relação a \texttt{PROJ}, a asserção abaixo \textbf{NÃO} expressa:



\begin{equation}\forall g \in \texttt{DSG}, \exists j \in \texttt{PROJ} : g.\texttt{PNO} = j.\texttt{PNO}\end{equation}
\begin{choices}
\choice Uma restrição a ser verificada na inserção de tuplas em $\texttt{DSG}$.
\choice Uma restrição que define um estado consistente do banco de dados.
\choice Uma restrição de integridade de chave estrangeira em $\texttt{DSG}$.
\choice Uma restrição a ser verificada na atualização de tuplas em $\texttt{DSG}$.
\choice Uma restrição de integridade de chave primária em $\texttt{PROJ}$.
\end{choices}

\question

O número de \textit{strings binárias} de comprimento 7 e contendo um par de zeros consecutivos é


\begin{choices}
\choice 92
\choice 91
\choice 94
\choice 90
\choice 95
\end{choices}

\question

Qual das seguintes afirmações sobre crescimento assintótico de funções não é verdadeira:
\begin{choices}
\choice \( 2^{n+1} = O(2^n) \)
\choice Se \( f(n) = O(g(n)) \) então \( g(n) = O(f(n)) \)
\choice Se \( f(n) = O(g(n)) \) e \( g(n) = O(h(n)) \) então \( f(n) = O(h(n)) \)
\choice \( 2n^2 + 3n + 1 = O(n^2) \)
\choice \( \log n^2 = O(\log n) \)
\end{choices}

\question

Quando Edgar Codd cricou o modelo relacional de dados,

também criou algumas regras para que a estrutura das tabelas fosse adequada às

operações de inserção, atualização, remoção e busca de dados. Ao avaliar um

banco de dados que armazena dados dos eleitores, dos partidos e dos candidatos

de uma eleição, você encontrou a situação ilustrada na tabela abaixo:

\begin{figure}[H]

  \centering

  \includegraphics[width=0.5\linewidth]{static/imgs/defaul.png}

\end{figure}

Em relação às boas normas de projeto do modelo relacional, assinale a

alternativa que corretamente identifica um problema no projeto desta tabela:


\begin{choices}
\choice A chave primária para a tabela não está definida e, somente olhando as colunas, não é possível dizer se é a coluna 'numero' ou a coluna 'titulo'.
\choice A tabela contém um atributo multivalorado (o correto seria que cada tabela não tivesse atributos multivalorados).
\choice Apesar do telefone ser um atributo multivalorado, do jeito que a tabela está criada só é possível armazenar dois números de telefones, criando uma restrição ao armazenamento de telefones de contado de eleitores que tenham mais de dois telefones.
\choice A tabela está misturando dados de três entidades diferentes, eleitores, partidos e candidatos (o correto seria que cada tabela armazenasse dados de um único assunto ou entidade).
\choice Todas as afirmativas anteriores.
\end{choices}

\question

Sobre a hierarquia de Chomsky podemos afirmar que:
\begin{choices}
\choice As linguagens reconhecidas por autômatos a pilha são as linguagens regulares
\choice Uma linguagem que é recursivamente enumerável não pode ser uma linguagem regular
\choice Há linguagens que não são nem livres de contexto nem sensíveis a contexto
\choice As linguagens livres de contexto e as linguagens sensíveis a contexto se excluem
\choice Uma linguagem que não é regular é livre de contexto
\end{choices}

\question

Em SQL, qual comando é utilizado para combinar linhas de duas ou mais tabelas?


\begin{choices}
\choice LINK
\choice MERGE
\choice COMBINE
\choice JOIN
\choice CONNECT
\end{choices}

\question

Na linguagem SQL, qual das seguintes palavras-chave é usada para modificar os

dados em uma tabela existente?


\begin{choices}
\choice INSERT
\choice Nenhuma das alternativas acima
\choice UPDATE
\choice SELECT
\choice CREATE
\end{choices}

\question

São motivos importantes e gerais para o estudo das linguagens de programação,

exceto:
\begin{choices}
\choice Melhorar o embasamento na escolha das linguagens para um projeto
\choice Aumentar a capacidade de expressar idéias
\choice Aprender a sintaxe e a semântica de uma determinada linguagem
\choice Entender a importância da implementação
\choice Avanço geral da computação
\end{choices}

\question

Todos os convidados presentes num jantar tomam chá ou café. Treze convidados bebem café, dez bebem chá e 4 bebem chá e café. Quantas pessoas tem nesse jantar.


\begin{choices}
\choice 19
\choice 27
\choice 23
\choice 10
\choice 15
\end{choices}

\question

Sejam os seguintes predicados de uma linguagem de primeira ordem:



\begin{itemize}

    \item \( N(x) \): \( x \) é número;

    \item \( P(x) \): \( x \) tem propriedade \( P \);

    \item \( x < y \): \( x \) é menor que \( y \).

\end{itemize}



E sejam os símbolos:

\begin{itemize}

    \item \( \forall \): quantificador universal;

    \item \( \Rightarrow \): operador se-então;

    \item \( \neg \): operador de negação.

\end{itemize}



Para a fórmula: \( \forall x (N(x) \Rightarrow \neg \forall y (N(y) \Rightarrow y < x)) \), qual alternativa abaixo \textbf{NÃO} constitui uma tradução possível?


\begin{choices}
\choice Para todo número, não é verdade que qualquer número seja menor do que ele.
\choice Para qualquer \( x \), se \( x \) é número, então não é verdade que todos os números são menores do que ele.
\choice Não há um número menor do que outro número.
\choice Para todo número, existe um outro número que é maior do que ele.
\choice Não há um número tal que todos os números são menores do que ele.
\end{choices}

\question

Seja a seguinte linguagem, onde \(\epsilon\) representa o string vazio e \$ representa um marcador de fim de entrada:



$S \rightarrow ABCD \\$\\

$A \rightarrow a \mid \epsilon \\$\\

$B \rightarrow a \mid \epsilon \\$\\

$C \rightarrow c \mid \epsilon \\$\\

$D \rightarrow S \mid c \mid \epsilon$\\



É incorreto afirmar que:
\begin{choices}
\choice O conjunto FOLLOW(A) = a, c, \$
\choice O conjunto FIRST(D) é igual ao conjunto FIRST(S)
\choice O conjunto FOLLOW(B) = c, \$
\choice O conjunto FOLLOW(D) é igual a FOLLOW(S)
\choice O conjunto FIRST(A) = a, \(\epsilon\)
\end{choices}

\question

Um banco de dados (BD) é uma coleção logicamente coerente de dados relacionados

e persistentes. Em relação às características de um banco de dados, assinale a

alternativa verdadeira:
\begin{choices}
\choice Geralmente é projetado e usado para finalidades gerais, sendo extremamente genéricos e adaptáveis para qualquer situação.
\choice Não precisa de dados atuais e verdadeiros para ser útil (para representar corretamente a realidade atual).
\choice Armazena dois grandes ``tipos`` de dados: os dados do usuário final e os metadados, sendo que os metadados ficam armazenados nas mesmas tabelas que os dados dos usuários e isso nos permite saber a estrutura dos bancos de dados.
\choice Não precisam ser obrigatoriamente computadorizados, existem bancos de dados manuais e/ou em papel.
\choice Representa um aspecto (ou parte) do mundo real, parte essa chamada de ``catálogo' ou ``dicionário de dados``.
\end{choices}

\question

Analise o diagrama relacional do banco de dados exibido na figura abaixo:

\begin{figure}[H]

  \centering

  \caption{Banco de dados de peças e fornecedores}

  \label{fig:pecas}

  %\fbox{

    \includegraphics[width=0.5\linewidth]{static/imgs/defaul.png}

  %}\\

  %\footnotesize{Fonte: xxxx.}

\end{figure}

Podemos dizer que esse banco de dados ilustra:
\begin{choices}
\choice Um relacionamento com cardinalidade N:N entre as tabelas ``pecas' e ``fornecedores', realizado através de relacionamentos identificados através da tabela ``fornecimentos'. 
\choice Um relacionamento com cardinalidade N:N entre as tabelas ``pecas' e ``fornecimentos', realizado através de relacionamentos não identificados com a tabela ``fornecedores'. 
\choice Um relacionamento com cardinalidade 1:N entre as tabelas ``pecas' e ``fornecedores', indicando que uma peça pode ter sido fornecida por diversos fornecedores. 
\choice Um relacionamento com cardinalidade N:N entre as tabelas ``pecas' e ``fornecedores', realizado através de relacionamentos não identificados através da tabela ``fornecimentos'. 
\choice Um relacionamento com identificação N:N entre as tabelas ``pecas' e ``fornecedores', realizado através do relacionamento com cardinalidade identificada através da tabela ``fornecimentos'. 
\end{choices}

\question

No que diz respeito as vantagens da arquitetura de micro-núcleo para sistemas operacionais em relação a arquiteturas de núcleo monolítico, quais das seguintes afirmações são verdadeiras?



\begin{itemize}

    \item[I] A arquitetura de micro-núcleo facilita a depuração do SO.

    \item[II] A arquitetura de micro-núcleo permite um número menor de mudanças de contexto.

    \item[III] A arquitetura de micro-núcleo facilita a reconfiguração de serviços do SO pois a maioria deles reside em espaço de usuário.

\end{itemize}
\begin{choices}
\choice Apenas I
\choice I e II
\choice II e III
\choice I e III
\choice Todas são verdadeiras
\end{choices}

\question

Você estava ouvindo a discussão de dois colegas, o João e a Maria. Enquanto o

João afirmava que o conceito de banco de dados (BD) é extremamente diferente do

conceito de sistema de gerenciamento de bancos de dados (SGBD), a Maria

discordava e afirmava que o banco de dados nada mais é do que um componente do

SGBD. Nesse momento o professor entrou na sala, percebeu a discussão e falou

que:
\begin{choices}
\choice O João estava certo já que o banco de dados é, realmente, um conceito diferente do conceito do sistema de gerenciamento de bancos de dados, e existe de forma independente.
\choice Os dois estavam certos pois o conceito, na prática, depende da situação específica que se está considerando.
\choice Os dois estavam errados pois não são conceitos separados nem componentes, são conceitos sinônimos.
\choice A Maria estava certa já que o banco de dados é, realmente, um componente interno do sistema de gerenciamento de bancos de dados e não existe de forma independente.
\choice Não é possível saber quem estava certo ou errado pois não definiram qual o modelo de dados que estava sendo considerado na discusão.
\end{choices}

\question

Ao analisar um novo sistema de gerenciamento de bancos de dados que surgiu no

mercado, você percebe que a entrada às operações do modelo de dados são tabelas

mas as saídas nunca são tabelas, são sempre outras estruturas. O vendedor desse

SGBD afirma que ele é um SGBD relacional. Pode-se afirmar que:


\begin{choices}
\choice Esse SGBD é relacional pois o vendedor afirmou que é.
\choice Nenhuma das alternativas anteriores.
\choice Esse SGBD não é relacional pois, para ser verdadeiramente relacional, as operações do modelo de dados não podem trabalhar recebendo tabelas como entrada.
\choice Esse SGBD não é relacional pois no modelo relacional a única estrutura percebida pelos usuários são tabelas.
\choice Esse SGBD é relacional pois consegue trabalhar com tabelas para a entrada das operações do modelo de dados. A saída não é importante.
\end{choices}



    %%%%%%%%%%%%%%%%%%%%%%%%%%%%%%%%%%%%%%%%%%%%%%%%%%%%%%%%%%%%%%%%%%%%%%%%%%%%%%%%
    %%% Finaliza o bloco de questões:
    \end{questions}
    \end{document}
    